\section{A full dyadic extension remembers its squares}\label{sec:extensions}
\subsection{The extension criterion with a fixed center}
Let $V$ be a finite-dimensional vector space over $\F_2$, and consider central extensions
\[
 1\longrightarrow Z\longrightarrow E_s\longrightarrow V\longrightarrow1,
 \qquad
 1\longrightarrow Z\longrightarrow E_t\longrightarrow W\longrightarrow1,
\]
where $Z$ is the same abelian group in both extensions. In applying this to a represented scalar pushout, the represented center and any already factored character kernel must be recorded. A pushout cannot recover invisible elements that have been discarded.

Fix a linear isomorphism $f:V\to W$. Choose an ordered basis $(e_i)$ of $V$, lifts $a_i$ in $E_s$, and lifts $b_i$ of $f(e_i)$ in $E_t$. Put $s_i=a_i^2$ and $t_i=b_i^2$ in $Z$. Write $Z^2=\{z^2:z\in Z\}$.

\begin{proposition}[Center-fixing extension criterion]\label{prop:ext}
There exists an isomorphism $F:E_s\to E_t$ fixing $Z$ pointwise and lifting $f$ if and only if
\[
 [a_i,a_j]=[b_i,b_j]\quad\hbox{for every }i,j,
 \qquad s_i t_i^{-1}\in Z^2\quad\hbox{for every }i.
 \tag{26.1}\label{eq:extcrit}
\]
If the generating lifts themselves must satisfy $F(a_i)=b_i$, then their squares must agree exactly, not merely modulo $Z^2$.
\end{proposition}
\begin{proof}
Any such isomorphism has $F(a_i)=\lambda_i b_i$ for some $\lambda_i\in Z$. Preservation of squares gives $s_i=\lambda_i^2 t_i$, while preservation of commutators gives the first condition. Conversely choose $\lambda_i$ with these squares. Every group element has a unique normal form
\[
 z a_1^{v_1}\cdots a_r^{v_r},\qquad z\in Z,\quad v_i\in\{0,1\}.
\]
Define its image by replacing $a_i$ with $\lambda_i b_i$. Normal-form multiplication uses only central elements, the displayed squares, and the pairwise commutators. The conditions ensure that each multiplication relation is preserved. The inverse construction proves bijectivity. Taking every $\lambda_i=1$ gives the final assertion.
\end{proof}

For $Z=\mu_2=\{1,-1\}$, each square is independent of the chosen lift and defines a quadratic refinement
\[
 a_v^2=(-1)^{q(v)},\qquad
 q(v+w)=q(v)+q(w)+B(v,w),
 \tag{26.2}\label{eq:quadratic}
\]
where $(-1)^{B(v,w)}$ is the commutator. A center-fixing extension isomorphism over a specified symplectic map $f$ therefore requires $q_t\circ f=q_s$.

For $Z=\mu_4$, multiplying a lift by $i$ changes its square by $-1$. The unmarked square classes are consequently weaker. For $Z=\C^\times$, every square ratio has a square root and the unmarked criterion reduces to matching commutators. Neither statement deletes a separately marked finite subgroup or a prescribed family of lifts.

\subsection{Dihedral and quaternion data inside the same scalar pushout}
Consider the matrices
\[
 X=\begin{pmatrix}0&1\\1&0\end{pmatrix},\qquad
 Z_0=\begin{pmatrix}1&0\\0&-1\end{pmatrix}.
\]
They satisfy $X^2=Z_0^2=1$ and $XZ_0=-Z_0X$. The group $H_D=\langle X,Z_0\rangle$ is a dihedral group of order eight. The group $H_Q=\langle iX,iZ_0\rangle$ is a quaternion group of order eight. On their common two-dimensional quotient, their refinements are
\[
 q_D(x,y)=xy,\qquad q_Q(x,y)=xy+x+y.
 \tag{26.3}\label{eq:DQ}
\]

\begin{center}\begin{tabular}{lrrrr}\toprule
Marked extension & $q(e)$ & $q(f)$ & $q(e+f)$ & $\sum_v(-1)^{q(v)}$\\\midrule
$H_D$ & 0 & 0 & 1 & $+2$\\
$H_Q$ & 1 & 1 & 1 & $-2$\\\bottomrule
\end{tabular}\end{center}

Both scalar pushouts identify with
\[
 P=\{zX^aZ_0^b:z\in\C^\times,\ a,b\in\{0,1\}\}.
\]
The marked subgroups remain different. In $H_D$, five nonidentity elements have order two and two have order four. In $H_Q$, only the central nonidentity element has order two and six have order four. In particular, no center-fixing automorphism of $P$ carries one marked subgroup onto the other.

On the other hand, a quaternion source does embed into $P$ through $iX,iZ_0$. Its squares and negative Gauss sign are retained. What fails is identifying it with the \emph{independently chosen} dihedral subgroup $H_D$. This example separates a successful transported marking from a false assertion about a fixed native marking.

\subsection{A full marked section contains more than a square table}
Suppose a normalized section $\sigma:V\to E$ has product cocycle
\[
 \sigma(v)\sigma(w)=c(v,w)\sigma(v+w).
\]
If an isomorphism must send each specified section element to a specified target section element, the two cocycles must agree after transport on \emph{every} pair. Equality of polar forms and square tables only identifies their extension class, not all those pointwise choices.

For example, start with the ordered section $\sigma(x,y)=X^xZ_0^y$ of $H_D$. Change only the value at $e+f$ by a central sign. Every square remains unchanged because $(-1)^2=1$, but the product of the specified lifts at $e$ and $f$ now differs by a sign from the specified lift at $e+f$. An isomorphism fixing the two generating lifts cannot also implement this changed third assignment. The new tests retain this counterexample rather than treating it as a harmless convention change.

\subsection{The valid scalar-pushout attachment contract}
There is nevertheless a positive attachment statement in the restricted sector already used by Paper XI. Suppose an actual type, its full character and its Heisenberg representation $\eta_W$ are supplied. Suppose the graph has one vertex, the prescribed stabilizer is trivial, and the source scalar pushout has the same faithfully represented center and a nondegenerate quotient $V$ with $\dim V\le\dim W$. No conflicting host-native finite marking is independently fixed.

Choose a symplectic embedding $f:V\to W$. Proposition~\ref{prop:ext} gives a scalar-center-fixing isomorphism of the source pushout with the restriction over $f(V)$. Carry every supplied finite subgroup, every lift, and every section value through this \emph{same} isomorphism. Since it is a homomorphism fixing scalars, it preserves all squares, commutators and section products, not just their classes.

Writing $W=f(V)\perp U$, the finite Heisenberg representation factors as $\eta_V\otimes\eta_U$. A projector onto a line in the second factor gives the source Heisenberg sector inside the actual image algebra. The inherited recovery theorem then returns the unchanged full inducing representation and candidate, under its stated recovery hypotheses. There is no stabilizer equation because the stipulated stabilizer is trivial.

This is an attachment into the scalar extension of the actual type. It is not a claim that every source lift is an unscaled arithmetic matrix, nor that the transported subgroup equals an independently specified canonical host subgroup. A package that records only the unmarked ambient scalar extension loses precisely the information displayed by \eqref{eq:DQ}.

\source{[Q16] contains the complete marked scalar-pushout contract. The normal-form proof, the dihedral/quaternion calculations and the section counterexample are explicit finite arguments. Their application uses the already supplied actual type and inherited recovery theorem. This section does not construct those arithmetic inputs.}
